\subsection{课程边界与 Tokenizer worked example}

slide 79 明确列出本讲未覆盖的架构、推理、多模态、滥用、上下文、数据墙与法律问题，并指向 CS224N、CS324、CS336 等后续学习路径。slides 80--86 是 BPE 的逐步 reveal；中间六页只增加步骤文字，最终页包含完整流程，因此覆盖矩阵保留 slide 86。这个收尾提醒读者：tokenizer 看似基础，却决定序列长度、词表、跨语言效率和模型实际处理的离散单位。

\lecturefigure{slides-images/slide-079.png}{官方 slide 79：本讲未覆盖主题与后续课程}
\lecturefigure{slides-images/slide-086.png}{官方 slide 86：BPE tokenizer 的完整训练流程}

\begin{knowledgebox}{读图：BPE 如何从字符得到子词}
BPE 从字符或字节级 token 开始，统计相邻 token pair 的频率，把最高频 pair 合并成新 token，并重复直到达到目标词表大小。若语料中 ``t h e'' 高频出现，算法可能先合并 ``t+h''，再合并 ``th+e''。常见片段变成单 token，可缩短序列；罕见词仍能拆成更小单位，避免 word-level unknown token。实际 tokenizer 还要规定 normalization、空格、特殊 token 和 byte fallback。
\end{knowledgebox}

\begin{warningbox}{Tokenizer 不是语言中立的压缩器}
词表由训练语料频率决定。英文占比高时，英文词可能用更少 token 表示，而低资源语言、代码或特殊符号被切得更碎，导致相同字符数消耗更多上下文和推理费用。改变 tokenizer 还会改变 embedding/output 层和已有 checkpoint 的兼容性，因此通常在预训练前固定。课程把它放到附加页，是基础复习，不代表它对 Agent 成本和多语言能力不重要。
\end{warningbox}

至此，官方 deck 的教学节点形成完整闭环：tokenizer 定义输入单位，预训练从大规模语料学习预测分布，SFT 与 RL 改变行为目标，评估决定优化方向，系统工程决定这些实验能否在预算内运行。后续章节不再逐页复述，而是把这条证据链重组为 Agent 项目可直接使用的训练、评估、成本与上线方法。

